\chapter{Dérivation HMT du paramètre de Barbero--Immirzi et de l’aire minimale : incidence hexade--octade, canaux \(90/120\) et spectre de l’opérateur d’aire}
\label{chap:p3-final:51:barbero_area}
\section{Morphisme d’incidence hexade--octade et conservation des canaux orientés \(90/120\)}

La dodécade et l’alignement du
\cref{sec:w12-w24-elevacion} étant fixés, soit \(H\) une hexade dans les coordonnées HMT,
soit \(\ell(H)\subset D\) son image et soit
\(O_H=\iota_D(\ell(H))\) son unique relèvement. En marquant un point
et une paire de points de l’hexade résiduelle, on obtient
\[
F_6(H)=\ell(H)\times\binom{\ell(H)}2,
\qquad
|F_6|=6\binom62=90.
\]
Si le point marqué peut parcourir l’octade :
\[
F_8(H)=O_H\times\binom{\ell(H)}2,
\qquad
|F_8|=8\binom62=120.
\]
L’inclusion \(\ell(H)\hookrightarrow O_H\) induit le morphisme d’incidence
\[
 \mathcal I_{6\to8}:F_6(H)\hookrightarrow F_8(H).
\]
En normalisant la différence par le facteur
choisi \(24\binom62\), on obtient
\[
\boxed{
\frac{90-120}{24\binom62}
=\frac{6-8}{24}
=-\frac1{12}.}
\]
L’égalité est un défaut normalisé de l’interface combinatoire
\(6\to8\). Le facteur \(24\binom62\) fait partie de la définition de cette
normalisation ; l’inclusion à elle seule détermine la différence \(-30\).
L’opérateur \(\mathcal I_{6\to8}\) a pour domaine les configurations
marquées d’incidence ; ce n’est ni l’extracteur transversal de l’état
dodécaphasique ni la famille des moments angulaires.

% Construcción dodecafásica del funcional de retorno.
\section{Chaîne fondamentale dodécaphasique et fonctionnelle hexadique--octadique}

L’incidence active se construit avant toute réalisation gravitationnelle.
Soit
\[
 H_{\mathrm{act}}=\{1,2,4,5,7,8\},
 \qquad
 \mathcal P_2(H_{\mathrm{act}})
 =\{A\subset H_{\mathrm{act}}:|A|=2\}.
\]
Alors
\[
 |H_{\mathrm{act}}|=6,
 \qquad
 |\mathcal P_2(H_{\mathrm{act}})|=15,
\]
et les incidences antérieure et postérieure à l’extension ont pour cardinaux
\[
 \mathcal F_{90}=H_{\mathrm{act}}\times
 \mathcal P_2(H_{\mathrm{act}}),
 \qquad |\mathcal F_{90}|=90,
\]
\[
 \mathcal F_{120}=\{1,\ldots,8\}\times
 \mathcal P_2(H_{\mathrm{act}}),
 \qquad |\mathcal F_{120}|=120.
\]
Ces deux cardinaux sont des sorties de la règle de phase ; ce ne sont ni des angles ni
des paramètres reçus de la réalisation physique.

Notons \(C_{12}=\mathbb Z/12\mathbb Z\) le registre temporel du TPK et
\(\partial_{\rm ret}\) son unique couture de retour par tour. Dans le
module libre des événements orientés, on considère la chaîne fondamentale
\begin{equation}
 c_{12}^{+}
 =\sum_{j\in C_{12}}[j,\mathcal F_{90}]
  -[\partial_{\rm ret},\mathcal F_{120}].
 \label{intsuccpass:eq:c12-fundamental-chain}
\end{equation}
La somme contient exactement les douze secteurs de la roue ; le second
terme est unique parce qu’un tour possède une seule couture de clôture. Son signe
est le signe de frontière. L’involution bidirectionnelle envoie
\(c_{12}^{+}\) sur \(c_{12}^{-}=-c_{12}^{+}\) : elle change l’orientation de la
chaîne, mais n’altère ni ses multiplicités absolues ni ses types
d’incidence.

Pour \(m\in\{90,120\}\), soit
\[
 S_m(q)=\frac{q^m}{1-q^{3m}}.
\]
Le lecteur de retour \(\mathscr R\) est défini sur les générateurs de
\eqref{intsuccpass:eq:c12-fundamental-chain} par
\[
 \mathscr R([j,\mathcal F_{90}])=S_{90},
 \qquad
 \mathscr R([\partial_{\rm ret},\mathcal F_{120}])=S_{120}.
\]
La première égalité est constante sur \(C_{12}\) par covariance sous la
rotation de phase ; la seconde a un domaine distinct et n’agit que sur la
couture étendue. Par linéarité,
\begin{equation}
 \boxed{\mathscr R(c_{12}^{+})=12S_{90}-S_{120}.}
 \label{intsuccpass:eq:dodecaphase-forces-12minus1}
\end{equation}
Ainsi, les coefficients \(12:-1\) sont l’image de la chaîne fondamentale : douze
cellules de phase et une frontière orientée. Ils ne sont sélectionnés ni au moyen de la valeur
conventionnelle du paramètre de Barbero--Immirzi, ni au moyen d’une comparaison
décimale, ni en imposant au préalable le coefficient de pôle qui sera calculé
ci-après.

\begin{theorem}[Coefficient de pôle de la chaîne dodécaphasique]
La fonctionnelle \(F(q)=\mathscr R(c_{12}^{+})(q)
=12S_{90}(q)-S_{120}(q)\), déterminée par les douze secteurs
et leur couture orientée, a pour coefficient \(1/24\) dans
\((1-q)^{-1}\). Son résidu complexe en \(q=1\) est \(-1/24\).
Le coefficient de pôle est une sortie de la fonctionnelle déjà construite.
\end{theorem}

\begin{proof}
Pour une fonctionnelle ayant un pôle simple en \(q=1\), nous écrivons
\[
 p_1(F):=\lim_{q\uparrow1}(1-q)F(q).
\]
L’identité
\[
 1-q^{3m}=(1-q)\sum_{j=0}^{3m-1}q^j
\]
implique
\[
 p_1(S_m)
 =\lim_{q\uparrow1}\frac{q^m}{\sum_{j=0}^{3m-1}q^j}
 =\frac1{3m}.
\]
Par linéarité,
\[
 p_1\bigl(\mathscr R(c_{12}^{+})\bigr)
 =\frac{12}{270}-\frac1{360}
 =\frac{48-3}{1080}=\frac1{24}.
\]
Ces égalités s’effectuent dans \(\mathbb Q\). Le résidu complexe
est le coefficient de \((q-1)^{-1}\), de sorte que
\[
 \operatorname*{Res}_{q=1}\mathscr R(c_{12}^{+})(q)
 =-p_1\bigl(\mathscr R(c_{12}^{+})\bigr)=-\frac1{24}.
\]
\end{proof}

L’équation diophantienne
\[
 \frac a{270}-\frac b{360}=\frac1{24}
 \quad\Longleftrightarrow\quad
 4a-3b=45
\]
reste un certificat arithmétique ultérieur. À elle seule, elle admet la
famille \((a,b)=(12+3t,1+4t)\) ; la chaîne
\eqref{intsuccpass:eq:c12-fundamental-chain} élimine cette ambiguïté avant le calcul,
car modifier \((12,1)\) altère la multiplicité de l’orbite ou de sa couture
et ne représente donc plus un tour fondamental du TPK.

Avec les canaux conjugués \(q_\pm\) déjà engendrés par la branche angulaire, on
obtient enfin
\[
 \mathscr R(c_{12}^{+})(q_-)-
 \mathscr R(c_{12}^{+})(q_+)
 =12\Delta S_{90}-\Delta S_{120}.
\]
C’est l’entrée interne de la fonctionnelle hexadique--octadique. Son évaluation
comme paramètre de Barbero--Immirzi et son insertion dans l’opérateur d’aire sont
des publications ultérieures : elles ne participent pas à la sélection de
\(c_{12}^{+}\), de \(90/120\), de \(12:-1\) ni de \(1/24\).

\paragraph{Falsificateurs.}
La construction est réfutée si le recensement de phase ne donne pas douze secteurs, si
un tour contient plus d’une couture, si les incidences n’ont pas pour
cardinaux \(90\) et \(120\), si l’involution n’inverse pas la chaîne complète,
ou si le coefficient exact de \((1-q)^{-1}\) de son image diffère de \(1/24\).

\label{mdg:chap:barbero}

La transition d’une hexade à une octade combine trois structures de types
différents : une incidence finie, deux séries de Lambert associées aux
canaux \(q_\pm\) et une identification au paramètre réel de l’action de
Holst. Les trois s’articulent par l’incidence, la fonctionnelle angulaire et
le dictionnaire aréal vers la géométrie de la connexion
d’Ashtekar--Barbero développé dans cette même résidence.

\section{Cardinalités d’incidence 90 et 120}

Soit \(H\) un ensemble de six points et soit \(O\) un ensemble de huit points
contenant une copie distinguée de \(H\). Considérons
\[
 \mathcal F_6=H\times\binom H2,
 \qquad
 \mathcal F_8=O\times\binom H2.
\]
L’inclusion distinguée \(H\hookrightarrow O\) induit
\[
 \iota_{6,8}:\mathcal F_6\hookrightarrow\mathcal F_8,
 \qquad
 (h,\{h_1,h_2\})\longmapsto(h,\{h_1,h_2\}).
\]
Ce morphisme d’incidences est l’antécédent combinatoire des
cardinaux \(90/120\) ; ce n’est ni l’extracteur dodécaphasique ni une série angulaire.

\begin{theorem}[Cardinalités de l’inclusion]
\[
 |\mathcal F_6|=6\binom62=90,
 \qquad
 |\mathcal F_8|=8\binom62=120.
\]
La différence normalisée par vingt-quatre coordonnées est
\[
 \frac{90-120}{24\binom62}=-\frac1{12}.
\]
\end{theorem}
\begin{proof}
Chaque élément de la première coordonnée peut être combiné avec n’importe laquelle des
quinze paires de \(H\). Le principe multiplicatif donne les deux cardinalités ;
la dernière identité résulte de la substitution de \(\binom62=15\).
\end{proof}

\begin{figure}[htbp]
\centering
\begin{tikzpicture}[
  font=\small,
  setnode/.style={circle,draw,minimum size=6mm,inner sep=0pt,
    font=\scriptsize\bfseries},
  box/.style={draw,rounded corners=1.5mm,align=center,
    text width=5.0cm,inner xsep=6pt,inner ysep=5pt},
  arr/.style={-{Stealth[length=2.2mm]},line width=.7pt}]

  % Inclusión geométrica. Las flechas parten del contorno y no atraviesan nodos.
  \node[setnode,draw=hmtblue,fill=hmtlightblue] (h1) at (-2.3,1.25) {1};
  \node[setnode,draw=hmtblue,fill=hmtlightblue] (h2) at (-1.1,1.9) {2};
  \node[setnode,draw=hmtblue,fill=hmtlightblue] (h3) at (.1,1.25) {3};
  \node[setnode,draw=hmtblue,fill=hmtlightblue] (h4) at (.1,-.15) {4};
  \node[setnode,draw=hmtblue,fill=hmtlightblue] (h5) at (-1.1,-.8) {5};
  \node[setnode,draw=hmtblue,fill=hmtlightblue] (h6) at (-2.3,-.15) {6};
  \node[setnode,draw=hmtgreen,fill=hmtlightgreen] (o7) at (-3.2,.55) {7};
  \node[setnode,draw=hmtgreen,fill=hmtlightgreen] (o8) at (1.0,.55) {8};
  \node[draw=hmtblue,dashed,rounded corners=4mm,
    fit=(h1)(h2)(h3)(h4)(h5)(h6),inner sep=4mm,
    label={[hmtblue,font=\bfseries]above:{hexada \(H\), \(|H|=6\)}}] (H) {};
  \node[draw=hmtgreen,line width=.8pt,rounded corners=6mm,
    fit=(H)(o7)(o8),inner sep=5mm,
    label={[hmtgreen,font=\bfseries]below:{octada \(O\), \(|O|=8\), \(H\subset O\)}}] (O) {};

  \node[box,draw=hmtblue,fill=hmtlightblue] (f6) at (5.3,1.35)
    {Incidences définies sur \(H\)\\[1mm]
     \(\mathcal F_6=H\times\binom H2\),\quad
     \(|\mathcal F_6|=6\binom62=90\)};
  \node[box,draw=hmtgreen,fill=hmtlightgreen] (f8) at (5.3,-1.0)
    {Extension des incidences à \(O\)\\[1mm]
     \(\mathcal F_8=O\times\binom H2\),\quad
     \(|\mathcal F_8|=8\binom62=120\)};
  \draw[arr,hmtblue] (O.east) -- ++(.45,0) |- (f6.west);
  \draw[arr,hmtgreen] (O.east) -- ++(.45,0) |- (f8.west);
  \draw[arr,hmtgold] (f6.south) -- node[right,align=left,font=\footnotesize]
    {inclusion \(\iota_{6,8}\)} (f8.north);

  \node[box,draw=hmtgold,fill=hmtlightgold,text width=10.9cm]
    (def) at (1.45,-3.80)
    {Défaut combinatoire orienté\\[-.2mm]
     avec la normalisation déclarée :\\[1mm]
     \(\displaystyle
       \frac{|\mathcal F_6|-|\mathcal F_8|}{24\binom62}
       =\frac{90-120}{24\cdot15}=-\frac1{12}.\)};
  \draw[arr,hmtgold] (f8.south) -- (f8.south |- def.north);

  \node[align=center,text=hmtgray,font=\footnotesize,text width=11.2cm]
    at (1.45,-5.50)
    {Cette identité appartient à l’application d’incidences. Elle n’identifie pas
     l’extracteur dodécaphasique, le semi-groupe de tours ni la fonctionnelle angulaire.};
\end{tikzpicture}

\caption{Inclusion d’un sous-ensemble distingué de six éléments dans un ensemble de
huit éléments et dénombrements d’incidence. Les cardinaux \(90\) et \(120\) ne fusionnent ni l’extracteur
dodécaphasique, ni la fonctionnelle angulaire, ni le semi-groupe des tours.}
\label{mdg:fig:incidencia-seis-ocho-rev9}
\end{figure}

Les indices initiaux de la hiérarchie angulaire sont ainsi associés à deux
espaces d’incidence : 90 configurations sur l’hexade et 120 sur
l’octade ambiante. Le quotient \(-1/12\) est une identité combinatoire de cette
normalisation.

\section{Séries de Lambert associées aux 2 secteurs d’incidence}

Pour \(0<q<1\), nous définissons simultanément
\[
 S_{90}(q)=\frac{q^{90}}{1-q^{270}},
 \qquad
 S_{120}(q)=\frac{q^{120}}{1-q^{360}},
 \qquad
 \Delta S_m=S_m(q_-)-S_m(q_+).
\]
Les développements géométriques
\[
 \Delta S_{90}=\sum_{j\geq0}s_{90+270j},
 \qquad
 \Delta S_{120}=\sum_{j\geq0}s_{120+360j}
\]
situent les deux termes dans le semi-groupe global du
\cref{mdg:chap:torres}.

La fonctionnelle angulaire associée à la transition \(6\to8\) est
\[
 K_{6\to8}=12\Delta S_{90}-\Delta S_{120},
\]
et sa précédence a été fixée par la chaîne fondamentale dodécaphasique :
les douze secteurs et l’unique couture orientée produisent d’abord le poids
\(12:-1\). Le coefficient de \((1-q)^{-1}\), égal à \(1/24\), est ensuite calculé comme sortie exacte ;
l’équation diophantienne qui suit certifie la paire déjà engendrée et ne la
sélectionne pas.
La coordonnée angulaire résultante est
\[
\Gamma_{6\to8}
=\frac{\pi A\Cstar}{180}+K_{6\to8}.
\]
Au point HMT,
\[
\begin{aligned}
 \Delta S_{90}&=2.95169439297457621139847987861\times10^{-4},\\
 \Delta S_{120}&=1.96835112502951720093738706217\times10^{-5},\\
 \frac{\pi A\Cstar}{180}&=0.270485322934140078465586458790\ldots,
\end{aligned}
\]
et, par conséquent,
\[
 \boxed{\Gamma_{6\to8}
 =0.27400767269445927474725526077398041940\ldots.}
\]
Les tableaux calculés avec la coordonnée angulaire conjuguée \(C^*\) arrondie à quinze décimales donnent
\(0.2740076726944593\) en arithmétique à double précision ; la valeur précédente
est l’évaluation multiprécision de la phase régionale complète.

\section{Certificat diophantien ultérieur des poids}

L’identité
\[
 \lim_{q\to1^-}(1-q)\frac{q^m}{1-q^n}=\frac1n
\]
associe à la combinaison \(aS_{90}-bS_{120}\), avant son évaluation
dans les canaux conjugués, le coefficient de \((1-q)^{-1}\)
\[
 \frac a{270}-\frac b{360}.
\]

\begin{theorem}[Poids entiers de coefficient de pôle \texorpdfstring{\(1/24\)}{1/24}]
Les solutions entières positives de
\[
 \frac a{270}-\frac b{360}=\frac1{24}
\]
sont
\[
 (a,b)=(12+3t,1+4t),\qquad t\in\ZZ_{\geq0}.
\]
La paire \((12,1)\) est l’unique solution minimale coordonnée par coordonnée et
l’unique de norme \(\ell^1\) minimale.
\end{theorem}
\begin{proof}
La multiplication par \(1080\) produit \(4a-3b=45\). Modulo trois, on obtient
\(a\equiv0\pmod3\) ; en écrivant \(a=3u\), il en résulte \(4u-b=15\). La
positivité implique \(u=4+t\), \(b=1+4t\), avec \(t\geq0\). Par conséquent,
\(a+b=13+7t\), ce qui démontre la minimalité.
\end{proof}

La primitivité \(\gcd(a,b)=1\) ne sélectionne pas à elle seule la paire minimale, car
il existe d’autres solutions primitives. Le tour fondamental du TPK a déjà
produit le poids \(12:-1\) ; le coefficient de pôle \(1/24\), la positivité et la
minimalité forment son certificat arithmétique ultérieur. Aucun de ces
objets ne se déduit de l’évaluation décimale de \(\Gamma_{6\to8}\).

\newpage
\section{Monotonie de la fonctionnelle par rapport à la coordonnée angulaire conjuguée}

Para \(m\in\{90,120\}\),
\[
 S_m'(q)=\frac{m q^{m-1}(1+2q^{3m})}{(1-q^{3m})^2}>0.
\]
De plus,
\[
 \frac{\dd q_-}{\dd C}=\frac\pi{180}q_-,
 \qquad
 \frac{\dd q_+}{\dd C}=-\frac\pi{180}q_+.
\]

\begin{proposition}[Monotonie stricte]
Pour \(1^\circ\leq C\leq3^\circ\), la fonction
\(C\mapsto\Gamma_{6\to8}(A,C)\) est strictement croissante. Par conséquent,
toute valeur de son image possède au plus un antécédent dans cet intervalle.
\end{proposition}
\begin{proof}
La dérivée de chaque différence \(\Delta S_m\) est positive. Dans l’intervalle
indiqué, on obtient
\(\Delta S'_{120}<5.35\times10^{-4}\), tandis que
\(\pi A/180>0.1273\). Par conséquent,
\(\Gamma'(C)>0.1267\). En \(C=\Cstar\),
\[
 \Gamma'(\Cstar)=0.1328995105618907\ldots.
\]
\end{proof}

L’injectivité locale exclut une seconde coordonnée angulaire proche donnant la
même évaluation de la fonctionnelle. Ce contrôle ne construit ni ne sélectionne
\(C^*\) : la coordonnée angulaire conjuguée a déjà été obtenue dans la branche régionale précédente.
La proposition quantifie uniquement la sensibilité de l’évaluation descendante
une fois \(A\) et \(C^*\) fixés ; elle n’autorise pas à inverser
\(\Gamma_{6\to8}\) ni une valeur physique de \(\gamma_{\mathrm{BI}}\) pour
les redéfinir.

\label{ch:modular}

\section{Inclusion hexade--octade et réalisations fonctionnelles associées}

Fixons un ensemble \(H\) de six points et un ensemble \(O\) de huit
points avec une inclusion distinguée \(j:H\hookrightarrow O\). Dans la
résidence de Witt, \(H\) est une hexade d'une dodécade et \(O\) son octade
relevée ; dans ce chapitre, on utilise seulement l'inclusion déjà fixée, et non la
valeur terminale de Barbero--Immirzi. Nous écrivons
\begin{equation}
\label{eq:modular-pairs-H}
\binom H2:=\{P\subset H:|P|=2\},
\qquad \left|\binom H2\right|=\binom62=15,
\end{equation}
et définissons deux espaces d'incidences de types distincts :
\begin{equation}
\label{eq:modular-F6-F8}
\mathcal F_6:=H\times\binom H2,
\qquad
\mathcal F_8:=O\times\binom H2.
\end{equation}
L'inclusion d'incidences est
\begin{equation}
\label{eq:modular-iota68}
\iota_{6,8}:\mathcal F_6\hookrightarrow\mathcal F_8,
\qquad
(h,P)\longmapsto(j(h),P).
\end{equation}
Elle est injective parce que \(j\) l'est et conserve la paire marquée \(P\).

\begin{theorem}[Décomptes d'incidence et défaut orienté]
\label{thm:modular-incidence-90-120}
Les deux espaces de \cref{eq:modular-F6-F8} satisfont
\begin{equation}
\label{eq:modular-counts-90-120}
|\mathcal F_6|=6\binom62=90,
\qquad
|\mathcal F_8|=8\binom62=120,
\qquad
|\mathcal F_8\setminus\iota_{6,8}(\mathcal F_6)|=30.
\end{equation}
La normalisation par les vingt-quatre coordonnées du couple
dodécade--vingt-quatre et par les quinze paires de \(H\) produit
\begin{equation}
\label{eq:modular-incidence-defect}
\boxed{
\delta_{6,8}^{\rm inc}
=\frac{|\mathcal F_6|-|\mathcal F_8|}
{24\binom62}=-\frac1{12}.}
\end{equation}
La même fraction s'obtient par la fonctionnelle réciproque associée
à la mémoire incorporée,
\begin{equation}
\label{eq:modular-reciprocal-defect}
\boxed{
\delta_{6,8}^{\rm rec}
=\frac{30}{120}-\frac{30}{90}=-\frac1{12}.}
\end{equation}
Ce sont deux réalisations exactes du même défaut orienté, mais elles ont
des domaines distincts et ne sont pas identifiées à \(\zeta(-1)\) sans un
entrelaceur supplémentaire.
\end{theorem}

\begin{proof}
Le principe multiplicatif appliqué à
\cref{eq:modular-F6-F8} donne \(6\cdot15=90\) et \(8\cdot15=120\).
L'image de \(\iota_{6,8}\) a quatre-vingt-dix éléments, de sorte que le
complémentaire en a trente. Substituer dans
\cref{eq:modular-incidence-defect} donne
\(-30/(24\cdot15)=-1/12\), tandis que
\cref{eq:modular-reciprocal-defect} est
\(1/4-1/3=-1/12\).
\end{proof}

La transition hexade--octade admet deux réalisations fonctionnelles qui ne
se déduisent pas l'une de l'autre. Le diagramme algébrique, formulé sur les espaces
qui justifient les décomptes, est
\begin{equation}
\label{eq:modular-68-two-readers}
\begin{aligned}
(H\xhookrightarrow{\,j\,}O)
&\longmapsto
\left(|\mathcal F_6|,|\mathcal F_8|\right)=(90,120)
\longmapsto\Gamma_{6\to8}=\gammaH,\\
(h_6,o_8,t)=(6,8,3)
&\longmapsto-\frac{o_8-h_6}{t}=-\frac23,
\end{aligned}
\end{equation}
où \(-2/3\) est le coefficient de la troisième coordonnée de la ligne
\(13\) du constructeur CKM. La première réalisation transforme
l'incidence en aire et en mémoire modulaire ; la seconde transforme la même
transition en paramètres de mélange angulaire.
On ne postule pas \(\boldsymbol\theta_{\rm CKM}=f(\gammaH)\) : les deux sorties
conservent un antécédent commun.

Le morphisme qui détermine les secteurs \(90/120\) est
\(\iota_{6,8}\). La grandeur \(\gammaH\) a déjà été déterminée en interne
par la construction hexade--octade et est réutilisée ici sans être reconstruite à
partir de sa représentation décimale. Les
conséquences opératorielles ne commencent qu'après fixation de l'échelle
élémentaire \(a_0\), qui est définie ci-dessous sans utiliser l'entropie ni le
produit terminal \(\gammaH a_0\).

\section{Correspondance entre l'intérieur nonadique et les degrés de liberté de bord}

Soit \(\mathbb T_3=\{0,1,2\}\). Tout \(x\in\mathbb Z/9\mathbb Z\) possède
une écriture en chiffres unique
\(x=x_0+3x_1\), avec \(x_i\in\mathbb T_3\), et tout
\(y\in\mathbb Z/27\mathbb Z\) une écriture unique
\(y=y_0+3y_1+9y_2\). C'est pourquoi
\begin{equation}
\label{eq:modular-bulk-boundary-729}
\mathcal B=(\mathbb Z/9\mathbb Z)^3,
\qquad
\mathcal B_\partial=(\mathbb Z/27\mathbb Z)^2,
\qquad
|\mathcal B|=|\mathcal B_\partial|=3^6=729.
\end{equation}
L'égalité des cardinaux s'accompagne de l'application explicite. Si
\(x_j=r_{2j-1}+3r_{2j}\) pour \(j=1,2,3\), nous définissons
\begin{equation}
\label{eq:modular-beta729}
\beta_{729}(x_1,x_2,x_3)
=\bigl(r_1+3r_2+9r_3,\;r_4+3r_5+9r_6\bigr).
\end{equation}

\begin{proposition}[Bijection généalogique \(729\leftrightarrow729\)]
\label{prop:modular-beta729}
L'application \(\beta_{729}:\mathcal B\to\mathcal B_\partial\) est bijective et
préserve le mot ordonné de six trits. Ce n'est pas, en général, un
isomorphisme de groupes additifs : regrouper les chiffres change l'endroit où
la retenue est enregistrée.
\end{proposition}

\begin{proof}
Extraire les six chiffres de trois coordonnées nonadiques et les regrouper en
deux triplets est une permutation des coordonnées de \(\mathbb T_3^6\). La
décomposition positionnelle est unique des deux côtés, de sorte que
l'opération possède un inverse. La mise en garde additive se vérifie, par exemple,
en additionnant des mots qui engendrent une retenue entre \(r_3\) et \(r_4\) : cette
frontière sépare les deux coordonnées d'ordre vingt-sept, mais traverse
deux coordonnées distinctes dans le regroupement nonadique.
\end{proof}

Sur le tore de bord \(\mathcal B_\partial\), nous prenons le bloc de sommets
\begin{equation}
\label{eq:modular-block-A}
A=\{0,1,\ldots,8\}\times\{0,1,\ldots,8\}.
\end{equation}
Ses liens orientés qui traversent le bord se séparent par direction :
\begin{align}
\partial A_{90}
&=\{((-1,j),(0,j)),((8,j),(9,j)):0\leq j\leq8\},
\label{eq:modular-boundary90}\\
\partial A_{120}
&=\{((i,-1),(i,0)),((i,8),(i,9)):0\leq i\leq8\},
\label{eq:modular-boundary120}
\end{align}
où toutes les coordonnées se lisent modulo vingt-sept. Ainsi
\begin{equation}
\label{eq:modular-boundary-counts}
\partial A=\partial A_{90}\sqcup\partial A_{120},
\qquad
|\partial A_{90}|=|\partial A_{120}|=18.
\end{equation}
Les indices \(90,120\) étiquettent la provenance du canal ; ils ne disent pas que
le bord possède quatre-vingt-dix ou cent vingt liens.

\begin{theorem}[Loi aréale de la coupe triadique]
\label{thm:modular-triadic-cut}
Dans le graphe cubique \(N\times N\times N\) à capacité unitaire, toute
coupe entre deux faces opposées possède une capacité d'au moins \(N^2\), et une
couche transversale atteint cette borne. Pour \(N=9\),
\begin{equation}
\label{eq:modular-mincut-entropy}
\mathcal A_{\min}=81,
\qquad
S_{\rm tri}=81\log3.
\end{equation}
\end{theorem}

\begin{proof}
Il existe \(N^2\) droites d'arêtes parallèles à la direction normale, deux à deux
disjointes. Toute coupe doit intercepter chacune d'elles, de sorte que sa capacité
est d'au moins \(N^2\). Une seule couche transversale contient exactement
\(N^2\) arêtes et réalise le minimum. Si chaque lien coupé porte un trit
maximalement mélangé, son entropie est \(\log3\) ; l'additivité sur les
quatre-vingt-un liens donne \cref{eq:modular-mincut-entropy}.
\end{proof}

L'entropie \(S_{\rm tri}\) mesure la capacité de la coupe cubique.
L'entropie modulaire des trente-six liens de
\cref{eq:modular-boundary-counts} mesure un autre état et ne doit pas lui être
égalisée.

\section{Entrelaceur collectif d'incidence entre intérieur et bord}

Dans la somme
\(\ell^2(\mathcal F_6)\oplus\ell^2(\mathcal F_8)\), nous définissons
\begin{equation}
\label{eq:modular-u90-u120}
u_{90}=\frac1{\sqrt{90}}\sum_{f\in\mathcal F_6}(\delta_f,0),
\qquad
u_{120}=\frac1{\sqrt{120}}\sum_{f\in\mathcal F_8}(0,\delta_f).
\end{equation}
Dans
\(\ell^2(\partial A_{90})\oplus\ell^2(\partial A_{120})\), soient
\begin{equation}
\label{eq:modular-v90-v120}
v_{90}=\frac1{\sqrt{18}}\sum_{e\in\partial A_{90}}(\delta_e,0),
\qquad
v_{120}=\frac1{\sqrt{18}}\sum_{e\in\partial A_{120}}(0,\delta_e).
\end{equation}
Les couples \((u_{90},u_{120})\) et \((v_{90},v_{120})\) sont orthonormaux.
Il existe donc une unique isométrie entre les sous-espaces collectifs
qui satisfait
\begin{equation}
\label{eq:modular-J68}
\mathcal J_{6\to8}u_{90}=v_{90},
\qquad
\mathcal J_{6\to8}u_{120}=v_{120}.
\end{equation}
Définissons
\begin{align}
D_{\rm inc}
&=90|u_{90}\rangle\langle u_{90}|
+120|u_{120}\rangle\langle u_{120}|,
\label{eq:modular-Dinc}\\
D_{\rm area}
&=90|v_{90}\rangle\langle v_{90}|
+120|v_{120}\rangle\langle v_{120}|.
\label{eq:modular-Darea}
\end{align}

\begin{theorem}[Transport exact des deux étiquettes]
\label{thm:modular-J68}
L'application \(\mathcal J_{6\to8}\) est unitaire entre les deux plans
collectifs et entrelace les opérateurs d'incidence et d'aire :
\begin{equation}
\label{eq:modular-J-intertwining}
\boxed{
\mathcal J_{6\to8}D_{\rm inc}
=D_{\rm area}\mathcal J_{6\to8}.}
\end{equation}
Par conséquent, toute combinaison polynomiale des étiquettes \(90,120\), en
particulier la combinaison primitive \(12:-1\), est transportée avec les
mêmes coefficients.
\end{theorem}

\begin{proof}
L'orthonormalité montre que \(\mathcal J_{6\to8}\) envoie une base
orthonormale sur une autre. En \(u_{90}\), les deux membres de
\cref{eq:modular-J-intertwining} valent \(90v_{90}\) ; en \(u_{120}\),
ils valent tous deux \(120v_{120}\). L'égalité sur une base prouve l'identité.
Le calcul fonctionnel polynomial conserve l'entrelacement.
\end{proof}

Cette application ne prétend pas injecter quatre-vingt-dix incidences individuelles dans
dix-huit liens : elle transporte exactement les deux modes uniformes et leurs
étiquettes spectrales. L'étendre à des fluctuations non uniformes exigerait
une nouvelle application d'incidence ; l'opérateur modulaire défini ici n'utilise que
le sous-espace collectif qui vient d'être construit.

\section{Détermination de l'échelle élémentaire du spectre de l'opérateur d'aire}

Soit \(C^*\) la coordonnée angulaire conjuguée, exprimée en degrés, et fixons
\begin{equation}
\label{eq:modular-theta-clock-a0}
\theta_{\rm clk}=\frac{C^*}{6}\frac{\pi}{180},
\qquad
\boxed{
a_0=\frac{1+2\cos(10^\circ)}{3}\,\theta_{\rm clk}.}
\end{equation}
Le facteur qui précède \(\theta_{\rm clk}\) n'est pas une correction
décimale. C'est le caractère réel normalisé du triplet de phases
\begin{equation}
\label{eq:modular-c10-character}
\frac13\Tr\operatorname{diag}
\left(1,\ee^{\ii\pi/18},\ee^{-\ii\pi/18}\right)
=\frac{1+2\cos(10^\circ)}3.
\end{equation}
Ainsi, \(a_0\) est la moyenne spectrale de la structure angulaire sur le trit
orienté. Elle est positive, est fixée avant l'opérateur d'aire et n'est pas
sélectionnée à partir de l'entropie. L'identification ultérieure
\(\gammaH=\Gamma_{6\to8}\) réalise la sortie HMT dans la connexion
d'Ashtekar--Barbero \cite{Barbero,Immirzi} ; elle ne rouvre pas sa dérivation.

\begin{remark}[Portée du symbole \(a_0\)]
Dans ce chapitre, \(a_0\equiv a_0^{\rm area}\) désigne exclusivement la
normalisation de \cref{eq:modular-theta-clock-a0}. Ce n'est ni le rayon de Bohr
ni le facteur d'échelle cosmologique initial, bien que ces objets s'écrivent
traditionnellement avec le même symbole. Aucune identité de ce chapitre
ne permet de substituer l'un à l'autre.
\end{remark}

\section{Loi aréale de la coupe triadique : capacité \(81\log 3\), opérateur d’aire, spectre et entrelaceur collectif des canaux \(90/120\)}
\label{hap:sec:area-entrelazamiento-barbero-rev11}
\index{Barbero--Immirzi!opérateur d’aire}
\index{entropie modulaire!Hamiltonien modulaire}

L’identification HMT de la fonctionnelle \(\Gamma_{6\to8}\) à la coordonnée
de Barbero--Immirzi s’effectue sur le même bord triadique qui porte
l’incidence (90/120). L’intérieur et son feuillet de bord ont la même cardinalité,
\[
 \mathcal B=(\mathbb Z/9\mathbb Z)^3,
 \quad |\mathcal B|=729,
 \qquad
 \partial\mathcal B=(\mathbb Z/27\mathbb Z)^2,
 \quad |\partial\mathcal B|=729,
\]
par le regroupement de six trits en trois coordonnées nonadiques ou deux
coordonnées d’ordre vingt-sept.

\begin{theorem}[Loi aréale de la coupe triadique]
\label{hap:thm:ley-areal-corte-triadico-rev11}
Dans le graphe cubique \(N\times N\times N\) de capacité unitaire, la coupe
minimale entre deux faces opposées a pour capacité \(N^2\). Pour \(N=9\),
\[
 \mathcal A_{\min}=81,
 \qquad S_{\rm tri}=\mathcal A_{\min}\log3=81\log3.
\]
\end{theorem}
\begin{proof}
Les \(N^2\) droites parallèles à la direction normale sont des chemins disjoints en
arêtes, de sorte que toute coupe les intercepte. Une couche transversale contient
exactement \(N^2\) arêtes et atteint la borne. Chaque liaison triadique apporte
\(\log3\) à l’entropie d’un état maximalement mélangé.
\end{proof}

Dans le tore de bord \(27\times27\), soit \(A\) le bloc canonique
\(9\times9\). Sa frontière se décompose exactement en
\begin{equation}
 \partial A=\partial A_{90}\sqcup\partial A_{120},
 \qquad |\partial A_{90}|=|\partial A_{120}|=18.
 \label{hap:eq:corte-90-120-rev11}
\end{equation}
Les liaisons horizontales représentent le canal (90), et les verticales le
canal (120). Cette décomposition transforme la paire d’incidences en une
décomposition opératorielle du bord.

L’entropie aréale \(S_{\rm tri}=81\log3\) appartient à l’état maximalement
mélangé des liaisons de la coupe. L’entropie bicouche
\(S_{\rm bi}(y)\) de la \cref{hap:def:entropia-bicapa-rev6} appartient à la
distribution normalisée entre les deux feuillets ; les deux objets conservent
des domaines distincts.

Soit \(N=\operatorname{diag}(0,1,2)\) l’opérateur nombre d’une liaison tritique.
La réalisation aréale déclare l’échelle positive
\[
 \theta_{\rm clk}=\frac{C^*}{6}\frac{\pi}{180},
 \qquad
 a_0=\frac{1+2\cos(10^\circ)}{3}\,\theta_{\rm clk},
 \qquad \gamma_{\rm HMT}=\Gamma_{6\to8}.
\]
La formule de \(a_0\) constitue la structure de départ de cette
réalisation ; une fois cette échelle fixée, l’opérateur et son spectre se déduisent
exactement.

\begin{definition}[Opérateur d’aire HMT]
Sur le feuillet de bord
\(\mathcal H_{\partial}=\bigotimes_{e\in\partial A}\mathbb C^3\), on définit
\begin{equation}
 \widehat{\mathcal A}_{\rm HMT}
 =\gamma_{\rm HMT}a_0\sum_{e\in\partial A}N_e.
 \label{hap:eq:operador-area-hmt-rev11}
\end{equation}
Son spectre dans la coupe canonique est
\[
 \operatorname{Spec}\widehat{\mathcal A}_{\rm HMT}
 =\{n\gamma_{\rm HMT}a_0:0\leq n\leq72\},
\]
avec des multiplicités données par les coefficients de
\((1+z+z^2)^{36}\). Le prolongement par coupes compatibles produit la
famille protégée \(\mathcal A_n^{\rm prot}=n\gamma_{\rm HMT}a_0\).
\end{definition}

L’entrelaceur qui fixe la normalisation se construit dans le sous-espace des
canaux collectifs. Soient \(u_{90},u_{120}\) les vecteurs unitaires uniformes
des incidences hexadique et octadique, et soient \(v_{90},v_{120}\) les
vecteurs unitaires uniformes des deux ensembles de liaisons de
\eqref{hap:eq:corte-90-120-rev11}. L’application
\begin{equation}
 \begin{aligned}
 \mathcal J_{6\to8}:{}
 &\operatorname{span}\{u_{90},u_{120}\}
 \longrightarrow
 \operatorname{span}\{v_{90},v_{120}\},\\
 &\mathcal J_{6\to8}u_{90}=v_{90},
 \qquad
 \mathcal J_{6\to8}u_{120}=v_{120}
 \end{aligned}
 \label{hap:eq:entrelazador-area-barbero-rev11}
\end{equation}
est une isométrie entre ces sous-espaces collectifs. Si
\[
 D_{\rm inc}=90|u_{90}\rangle\langle u_{90}|
 +120|u_{120}\rangle\langle u_{120}|,
 \qquad
 D_{\rm area}=90|v_{90}\rangle\langle v_{90}|
 +120|v_{120}\rangle\langle v_{120}|,
\]
alors
\[
 \boxed{\mathcal J_{6\to8}D_{\rm inc}
 =D_{\rm area}\mathcal J_{6\to8}.}
\]
En particulier, la combinaison primitive \(12:-1\) agit
avec les mêmes coefficients avant et après \(\mathcal J_{6\to8}\). Cette
égalité fixe le dictionnaire de normalisation qui insère
\(\gamma_{\rm HMT}=\Gamma_{6\to8}\) dans la connexion
d’Ashtekar--Barbero.


Le passage de l’hexade à l’octade est une transformation de la structure d’incidence.\par

Lors du passage de \(6\) à \(8\), les paires disponibles, les relations internes et la manière dont l’orientation peut se conserver changent. Cette transition engendre un défaut orienté. La variation cardinale réalise une rotation du mode d’organisation.\par

La même inversion de feuillet s’exprime par\par

\[
C^*\longmapsto -C^*.
\]

Mais des lecteurs différents répondent de manières différentes à cette inversion.\par

La fonctionnelle hexade–octade qui conduit au paramètre de Barbero–Immirzi est impaire : lorsque le feuillet est inversé, elle change de signe. L’orientation géométrique s’inverse.\par

L’entropie bicouche, en revanche, est paire : elle attribue la même entropie à \(C^*\) et à \(-C^*\) lorsque l’on observe la distribution de probabilités. L’information totale peut rester identique même si l’orientation a été inversée.\par

L’espace CKM réalise une troisième réponse. L’inversion induit sur les trois angles une réflexion rationnelle exacte :\par

\[
R_{\rm CKM}
=
M_{\rm CKM}
\operatorname{diag}(1,-1,1)
M_{\rm CKM}^{-1},
\]

con\par

\[
R_{\rm CKM}^2=I,
\qquad
\det R_{\rm CKM}=-1.
\]

Cela signifie que le changement de feuillet est transporté dans l’espace de mélange comme une véritable réflexion. Il existe une direction impaire de rang un et un plan qui reste fixe.\par

Le contre-angle \(C^*\) est précisément la coordonnée qui change. La phase construite à partir de \(A\) reste invariante. La transformation sépare ainsi deux mémoires :\par

\begin{itemize}[leftmargin=*,itemsep=0.35em]
\item une mémoire d’orientation ;
\item une mémoire de phase.
\end{itemize}

Cela offre une lecture très profonde de CKM. Les angles de mélange sont les coordonnées d’un espace sur lequel agit la même involution qui organisait déjà la transition hexade–octade.\par

La rotation \(6\to8\) fournit l’opération d’orientation que CKM réalise dans son propre codomaine ; la causalité physique appartient au transducteur correspondant.\par

Nous avons alors trois publications sœurs du même passage :\par

\begin{enumerate}[leftmargin=*,itemsep=0.65em]
\item \textbf{Publication géométrique.}
Le paramètre de Barbero–Immirzi oriente le quantum d’aire.\par

\item \textbf{Publication informationnelle.}
L’état réduit et son hamiltonien modulaire mesurent quelle information de cette orientation demeure au bord.\par

\item \textbf{Publication de mélange.}
La réflexion CKM transforme les coordonnées de mélange et isole leur composante impaire.\par
\end{enumerate}

\(\gamma_{\rm BI}\), l’entropie et les angles CKM demeurent des réalisations distinctes d’une charnière généalogique commune. L’unité réside dans l’antécédent antérieur à la ramification. Ils partagent la charnière, mais chacun en conserve une partie différente.\par

C’est l’une des plus belles expressions de la méthodologie HMT : l’unité fondamentale est retrouvée en remontant les généalogies jusqu’à trouver l’opération commune qui engendre des résultats terminaux différents.\par

La gravité conserve la rotation comme orientation de l’aire.\par
L’information conserve la rotation comme mémoire de l’état.\par
Le mélange conserve la rotation comme réflexion des coordonnées.\par

La même transformation « respire » de trois manières.\par


\section{Séparation causale entre le spectre d’aire, la mémoire modulaire et la réalisation Cartan--Holst}
La dérivation de $\gamma_{\mathrm{BI}}$, de l’aire minimale et du spectre de l’opérateur d’aire s’achève dans ce chapitre. La reconstruction conjointe de l’occupation et de la provenance par le hamiltonien modulaire de bord, et la réalisation Palatini--Cartan--Holst, conservent leurs domaines, opérateurs et preuves propres dans la Partie IV.
